\documentclass[11pt]{article}

\usepackage{amsmath,amssymb,amsthm,mathtools}
\usepackage[T1]{fontenc}
\usepackage{lmodern}
\usepackage[margin=1in]{geometry}
\usepackage{enumitem}
\usepackage{hyperref}
\usepackage{microtype}
\usepackage{booktabs}

\hypersetup{
  colorlinks=true,
  linkcolor=blue,
  citecolor=blue,
  urlcolor=blue
}

\newtheorem{theorem}{Theorem}[section]
\newtheorem{proposition}[theorem]{Proposition}
\newtheorem{corollary}[theorem]{Corollary}
\newtheorem{definition}[theorem]{Definition}
\newtheorem{lemma}[theorem]{Lemma}
\newtheorem{remark}[theorem]{Remark}

\newcommand{\Z}{\mathbb Z}
\newcommand{\one}[1]{\mathbf 1_{\{#1\}}}

\title{A compact cleanup of the remaining 033/062 structural block}
\author{}
\date{March 2026}

\begin{document}
\maketitle

\begin{abstract}
This note handles the last localized structural block that still sat outside the
cleaned odd-$m$ D5 package after 092 and the reproof notes 095--097.
It does \emph{not} address the separate compact-bridge task around 076.

It treats 033 and 062 asymmetrically.
It rewrites 062 as one compact first-exit theorem with the pre-exit
patch-avoidance argument spelled out.
For 033, it compresses the needed content to the single exact trigger-family
lemma
\[
H_{L1}=\{(q,w,u,\Theta)=(m-1,m-1,u,2):u\neq 2\}.
\]
The note does \emph{not} attempt to rederive the full defect-template quotient
behind 033 from scratch.
Instead it isolates the exact theorem content needed later and records its
pilot-range artifact verification.

So 098 removes 062 as a separate support theorem and localizes the remaining
033 content to one exact trigger-family lemma inside the package.
\end{abstract}

\section{Scope and status}

Fix an odd modulus $m$ in the accepted D5 regime and the best-seed channel
\[
R1 \to H_{L1}
\]
with source seed
\[
L=[2,2,1], \qquad R=[1,4,4].
\]
We use the current coordinates
\[
(q,w,u,\Theta)
\]
along the active branch after the alternate-$2$ entry out of $R1$.

The exact downstream use of the old 033/062 block is narrow:
\begin{enumerate}[label=\textup{(\arabic*)},leftmargin=2.5em]
\item identify the exact trigger family $H_{L1}$;
\item derive the universal family-dependent first-exit targets;
\item prove that the actual active branch agrees with the candidate branch up
      to first exit.
\end{enumerate}

The point of the present note is that only item~(1) still belongs to the
broader defect-template classification. The rest can be proved cleanly inside
one short structural theorem. The separate compact-bridge task around 076 is
not addressed here.

\begin{remark}[honest dependency boundary]
A genuine first-principles reproof of the full 033 defect-family quotient would
be a larger project than the odd-$m$ D5 globalization package itself needs.
What the package uses is only the exact $H_{L1}$ trigger formula. This note
therefore rewrites 033 to that single lemma and fully rewrites 062.
\end{remark}

\section{The only 033 theorem content needed downstream}

\begin{proposition}[exact trigger family]\label{prop:trigger-family}
For the unresolved best-seed channel $R1\to H_{L1}$, the target hole family is
exactly
\[
H_{L1}=\{(q,w,u,\Theta)=(m-1,m-1,u,2):u\neq 2\}.
\]
Equivalently:
\begin{enumerate}[label=\textup{(\arabic*)},leftmargin=2.5em]
\item the target layer is $\Theta=2$;
\item the target current coordinates satisfy $q=w=m-1$;
\item the only excluded fiber is $u=2$.
\end{enumerate}
\end{proposition}

\begin{remark}[proof status]
Proposition~\ref{prop:trigger-family} is the promoted theorem content of 033.
This note does not rebuild the full defect-template classification from the raw
artifact; it isolates the exact formula that later theorems use.
\end{remark}

\begin{proposition}[pilot-range exact verification of the trigger formula]
The artifact file
\path{artifacts/d5_defect_splice_transducer_033/data/best_seed_defect_templates.json}
verifies the exact $H_{L1}$ signature on the saved pilot moduli $m=5,7,9$:
\begin{center}
\begin{tabular}{cccc}
\toprule
$m$ & target count & parameter values & representative $(q,w,u,\Theta)$ \\
\midrule
$5$ & $20=5\cdot 4$ & $\{0,1,3,4\}=\Z_5\setminus\{2\}$ & $(4,4,4,2)$ \\
$7$ & $42=7\cdot 6$ & $\{0,1,3,4,5,6\}=\Z_7\setminus\{2\}$ & $(6,6,4,2)$ \\
$9$ & $72=9\cdot 8$ & $\{0,1,3,4,5,6,7,8\}=\Z_9\setminus\{2\}$ & $(8,8,4,2)$ \\
\bottomrule
\end{tabular}
\end{center}
So the saved defect-family artifact matches exactly the formula in
Proposition~\ref{prop:trigger-family} on the pilot range.
\end{proposition}

\begin{proof}
This is a direct readout from the saved JSON. For each of $m=5,7,9$, the
artifact records:
\begin{enumerate}[label=\textup{(\arabic*)},leftmargin=2.5em]
\item target count $m(m-1)$,
\item parameter values $\Z_m\setminus\{2\}$,
\item representative coordinates $(q,w,\Theta)=(m-1,m-1,2)$.
\end{enumerate}
That is exactly the finite-range signature of
\[
H_{L1}=\{(m-1,m-1,u,2):u\neq 2\}.
\]
This is an exact verification on the saved pilot range, not a replacement for a
full general proof of the defect-template classification.
\end{proof}

\section{Candidate orbit and phase-$1$ invariant}

\begin{definition}[source residue and candidate orbit]
Let
\[
\rho:=u_{\mathrm{source}}+1 \pmod m,
\]
where $u_{\mathrm{source}}$ is the source-$u$ label of the chosen $R1$ family.
Start at the alternate-$2$ entry
\[
E=(m-1,1,u_{\mathrm{source}},2).
\]
Define the candidate orbit $\widehat x_n=(q_n,w_n,u_n,\Theta_n)$ by
\[
\widehat F(q,w,u,0)=(q+1,w,u,1),
\]
\[
\widehat F(q,w,u,1)=(q,w,u+1,2),
\]
\[
\widehat F(q,w,u,2)=\bigl(q,w+\one{q=m-1},u,3\bigr),
\]
\[
\widehat F(q,w,u,\Theta)=(q,w,u,\Theta+1)
\qquad (3\le \Theta\le m-1).
\]
\end{definition}

On the section
\[
\Sigma:=\{\Theta=2\},
\]
this candidate orbit induces the odometer return map
\[
(q,w)\mapsto(q+1,\;w+\one{q=m-1}).
\]
Equivalently, with
\[
\eta(q,w):=q+m(w-1)-(m-1)\in\{0,\dots,m^2-1\},
\]
one has $\eta\mapsto \eta+1$.

The only two phase-$1$ states whose direction-$2$ or direction-$1$ successor can
reach $(q,w)=(m-1,m-1)$ are
\[
A:=(m-1,m-2,1),
\qquad
B:=(m-2,m-1,1).
\]
Their odometer indices are
\[
\eta(A)=m(m-3),
\qquad
\eta(B)=m(m-2)-1.
\]
So $A$ is encountered before $B$.

\begin{lemma}[phase-$1$ source-residue invariant]\label{lem:phase1-rho}
Along the candidate orbit, every phase-$1$ current state satisfies
\[
q\equiv u-\rho+1 \pmod m.
\]
In particular,
\[
u(A)\equiv \rho-2 \pmod m,
\qquad
u(B)\equiv \rho-3 \pmod m,
\]
where $u(A)$ and $u(B)$ denote the $u$-coordinates at those states.
\end{lemma}

\begin{proof}
The first phase-$1$ state after entry is reached from
\[
(m-1,1,u_{\mathrm{source}},2)\to(m-1,2,u_{\mathrm{source}},3)
\to\cdots\to(m-1,2,u_{\mathrm{source}},0)\to(0,2,u_{\mathrm{source}},1).
\]
There one has
\[
q=0,
\qquad
u=u_{\mathrm{source}}=\rho-1,
\]
so the claimed congruence holds.

From one phase-$1$ state to the next, the phase-$1$ witness step increments
$u$ by $1$, the later phase-$0$ witness step increments $q$ by $1$, and no
other step changes either coordinate. Hence both sides of
\[
q\equiv u-\rho+1 \pmod m
\]
increase by $1$, so the congruence is preserved.
Substituting $q(A)=m-1$ and $q(B)=m-2$ gives the formulas for $u(A)$ and
$u(B)$.
\end{proof}

\begin{lemma}[phase-$1$ trigger criterion]\label{lem:trigger-criterion}
Let $x$ be a phase-$1$ current state on the candidate orbit.
An alternate step from $x$ lands in $H_{L1}$ if and only if one of the
following holds:
\begin{enumerate}[label=\textup{(\arabic*)},leftmargin=2.5em]
\item $x=A=(m-1,m-2,1)$ and $u(x)\neq 2$, in which case direction $2$ lands in
      $H_{L1}$;
\item $x=B=(m-2,m-1,1)$ and $u(x)\neq 2$, in which case direction $1$ lands in
      $H_{L1}$.
\end{enumerate}
\end{lemma}

\begin{proof}
By Proposition~\ref{prop:trigger-family}, the successor must satisfy
\[
(q',w',\Theta')=(m-1,m-1,2),
\qquad
u'\neq 2.
\]
Only phase-$1$ states can step to layer $2$, and only directions $1$ and $2$
change $q$ or $w$.
So the only two ways to reach $(q',w')=(m-1,m-1)$ are:
\begin{itemize}[leftmargin=2em]
\item direction $2$ from $(q,w)=(m-1,m-2)$,
\item direction $1$ from $(q,w)=(m-2,m-1)$.
\end{itemize}
Those are exactly $A$ and $B$.
Directions $1$ and $2$ preserve $u$, so the extra condition is simply
$u\neq 2$.
\end{proof}

\section{Universal first exits and pre-exit patch avoidance}

\begin{proposition}[candidate first exits]\label{prop:candidate-first-exits}
On the candidate orbit:
\begin{enumerate}[label=\textup{(\arabic*)},leftmargin=2.5em]
\item if $\rho\neq 4$ (regular family), the first exit to $H_{L1}$ occurs at
      $A=(m-1,m-2,1)$ by direction $2$;
\item if $\rho=4$ (exceptional family), the state $A$ does not exit and the
      first exit to $H_{L1}$ occurs at $B=(m-2,m-1,1)$ by direction $1$.
\end{enumerate}
Equivalently,
\[
T_{\mathrm{reg}}=(m-1,m-2,1),
\qquad
T_{\mathrm{exc}}=(m-2,m-1,1).
\]
\end{proposition}

\begin{proof}
By Lemma~\ref{lem:trigger-criterion}, only the two phase-$1$ states $A$ and $B$
can possibly trigger an exit to $H_{L1}$.
Because $\eta(A)<\eta(B)$, state $A$ is always encountered first.

At $A$, Lemma~\ref{lem:phase1-rho} gives
\[
 u(A)\equiv \rho-2 \pmod m.
\]
So $u(A)=2$ if and only if $\rho=4$.
Hence:
\begin{itemize}[leftmargin=2em]
\item if $\rho\neq 4$, then $u(A)\neq 2$, so Lemma~\ref{lem:trigger-criterion}
      gives an exit at $A$ by direction $2$, and because $A$ is encountered
      first, this is the first exit;
\item if $\rho=4$, then $u(A)=2$, so $A$ does not exit.
\end{itemize}

Now suppose $\rho=4$.
Then Lemma~\ref{lem:phase1-rho} gives
\[
 u(B)\equiv \rho-3\equiv 1 \pmod m,
\]
so $u(B)\neq 2$.
Lemma~\ref{lem:trigger-criterion} therefore gives an exit at $B$ by direction
$1$.
Since $A$ was the only earlier possible trigger and is blocked when $\rho=4$,
this is the first exit.
\end{proof}

\begin{lemma}[only patch-table facts used]\label{lem:patch-signatures}
For the pre-exit argument below, we use only the following coarse coordinate
signatures of the six modified current classes in the mixed witness patch table:
\begin{enumerate}[label=\textup{(\arabic*)},leftmargin=2.5em]
\item the classes $R1,R2,R3,L2$ require $w=0$;
\item the class $L3$ requires $w=1$;
\item a pre-exit occurrence of $L1$ would require the branch to have already
      reached the $w=m-1$ edge of the corridor.
\end{enumerate}
No finer information from the patch table is used.
\end{lemma}

\begin{lemma}[pre-exit patch avoidance on the candidate orbit]\label{lem:patch-avoidance}
Let $N_f$ be the first-exit time of the candidate orbit in the chosen source
family $f$.
For every $0\le n<N_f$, the candidate state $\widehat x_n$ avoids all patched
current classes.
\end{lemma}

\begin{proof}
The value $w$ starts at $1$ and changes only by the increment
\[
w\mapsto w+1
\]
at phase $2$ with $q=m-1$.
Hence $w$ never decreases, and $w=0$ never occurs before exit.
By Lemma~\ref{lem:patch-signatures}(1), the classes $R1,R2,R3,L2$ are therefore
impossible.

The value $w=1$ occurs only at the entry state $(m-1,1,u_{\mathrm{source}},2)$,
because the first phase-$2$ wrap immediately creates $w=2$ and thereafter $w$
never decreases.
So by Lemma~\ref{lem:patch-signatures}(2), the class $L3$ cannot occur after
entry.
Since the entry state itself is the alt-$2$ start state rather than a patched
class, $L3$ is absent on the whole pre-exit segment.

On a regular family, Proposition~\ref{prop:candidate-first-exits} says the
branch exits already at $A=(m-1,m-2,1)$, before $w$ can ever reach $m-1$.
So $L1$ is impossible there.

On the exceptional family, the first creation of $w=m-1$ occurs at phase $3$
from the state $(m-1,m-2,2)$.
The next phase-$1$ state then has $(q,w)=(0,m-1)$, not $(m-1,m-1)$.
The first later phase-$1$ state with $w=m-1$ that could also carry $q=m-1$
would occur strictly after $B=(m-2,m-1,1)$, but
Proposition~\ref{prop:candidate-first-exits} says the branch exits already at
$B$.
So $L1$ is impossible before exceptional exit as well.

Thus no patched current class is encountered before first exit.
\end{proof}

\begin{theorem}[actual first exits and pre-exit $B$-region invariance]
\label{thm:actual-first-exits}
For each source family:
\begin{enumerate}[label=\textup{(\arabic*)},leftmargin=2.5em]
\item the actual active branch agrees with the candidate orbit up to first
      exit;
\item every pre-exit actual current state is $B$-labeled;
\item the first exit to $H_{L1}$ occurs at the universal family-dependent
      target
      \[
      T_{\mathrm{reg}}=(m-1,m-2,1)
      \quad\text{or}\quad
      T_{\mathrm{exc}}=(m-2,m-1,1),
      \]
      with regular families exiting at $T_{\mathrm{reg}}$ by direction $2$ and
      the exceptional family exiting at $T_{\mathrm{exc}}$ by direction $1$.
\end{enumerate}
\end{theorem}

\begin{proof}
Fix a source family and let $N_f$ be the candidate first-exit time from
Proposition~\ref{prop:candidate-first-exits}.
We prove by induction on $n\le N_f$ that the actual state $x_n$ agrees with the
candidate state $\widehat x_n$ in the full current coordinates
$(q,w,u,\Theta)$.

At $n=0$, both states are the same entry state
\[
E=(m-1,1,u_{\mathrm{source}},2).
\]
By Lemma~\ref{lem:patch-avoidance}, this state is not patched, hence it is
$B$-labeled.

Now assume $x_n=\widehat x_n$ for some $n<N_f$.
Lemma~\ref{lem:patch-avoidance} says $\widehat x_n$ is not patched, so the
actual state $x_n$ is $B$-labeled.
Therefore its witness update is the unmodified mixed-witness update, which
induces exactly the candidate step on $(q,w,u,\Theta)$: direction $1$ changes
only $q$, direction $4$ changes only $u$, direction $2$ changes only $w$, and
all other anchors leave $(q,w,u)$ fixed while increasing $\Theta$ by $1$.
Hence
\[
(q,w,u,\Theta)(x_{n+1})=(q,w,u,\Theta)(\widehat x_{n+1}).
\]
So the actual and candidate branches agree for all $n\le N_f$, and every
pre-exit actual state is $B$-labeled.

For $n<N_f$, Proposition~\ref{prop:candidate-first-exits} and
Lemma~\ref{lem:trigger-criterion} say that the candidate state
$\widehat x_n$ does not exit to $H_{L1}$; because the trigger test depends only
on the current coordinates $(q,w,u,\Theta)$, the actual state $x_n$ does not
exit either.
At time $N_f$, Proposition~\ref{prop:candidate-first-exits} gives an exit from
$\widehat x_{N_f}$ at the stated universal target and direction, so the actual
branch exits there as well.
This proves all three claims.
\end{proof}

\begin{corollary}[what 098 removes]
The theorem-level content of 062 is now inlined: no separate 062 support theorem
is needed for the odd-$m$ D5 package.
The old 033/062 block now survives inside the package only as the exact
trigger-family lemma isolated in Proposition~\ref{prop:trigger-family} and the
broader defect-template provenance behind it.
\end{corollary}

\section{Exact $m=5$ reference slice}

\begin{proposition}[$m=5$ structural specialization]
At $m=5$,
\[
H_{L1}=\{(4,4,u,2):u\neq 2\},
\qquad
T_{\mathrm{reg}}=(4,3,1),
\qquad
T_{\mathrm{exc}}=(3,4,1).
\]
So the two possible first-exit states are exactly
\[
A=(4,3,1),
\qquad
B=(3,4,1),
\]
and the exceptional family is the source residue class $\rho=4$.
\end{proposition}

\begin{proof}
Substitute $m=5$ into Proposition~\ref{prop:trigger-family} and
Proposition~\ref{prop:candidate-first-exits}.
\end{proof}

\section{Bottom line}

After 098, the odd-$m$ D5 package no longer needs 062 as a separate imported
support theorem. Together with 095--097, this leaves the compact concrete
bridge package around 076 as the main remaining imported theorem layer.
What survives of the old 033/062 block inside the package is only the exact
trigger-family lemma now isolated in Proposition~\ref{prop:trigger-family} and
its broader defect-template provenance. It is no longer a structural or
dynamical bottleneck for the globalization proof chain.

\end{document}
